\documentclass[12pt]{article}

% The preamble file contains all the macros and page setup. I recommend making a folder where you compile your exercises and always call on the preamble.
%%%%%%%%%%%%%%%%%%%%%%%%%%%%%%%%%%%%%%%%%%%%%%%%%%%%%%%%%%%%%%%%%%%%%%
%% all the formatting stuff and packages



\usepackage{geometry}
\geometry{
  top=4cm,
  headheight=70pt,
  headsep=25pt,
  inner=1.25in,
  outer=1.25in
}

\usepackage[utf8]{inputenc}

\usepackage{xcolor,graphicx,hyperref,amsthm,amsmath,amssymb}
\usepackage{algpseudocode}


\usepackage{fancyhdr}
\pagestyle{fancy}

%%% Begin Header Definitions %%%
\fancypagestyle{firststyle}
{
   	\lhead{}
	\rhead{\raisebox{0.2cm}{\includegraphics[width=3.5cm]{resource/hpi_logo.pdf}}}
	\renewcommand{\headrulewidth}{0pt}
}


\newcommand{\makeheader}[1]{
\noindent
\textbf{Algorithmix \hfill Winter term 2019/20 \\ Andreas G\"obel}

\begin{center}
\large{\textbf{\makeHeadline}}\\
\end{center}

\begin{center}
\textit{#1} 
\end{center}


}

\newcommand{\lecturer}[1]{\newcommand{\makeLecturer}{#1}}
\newcommand{\lectureName}[1]{\newcommand{\makeLectureName}{#1}}
\newcommand{\semester}[1]{\newcommand{\makeSemester}{#1}}
\newcommand{\headline}[1]{\newcommand{\makeHeadline}{#1}}
\newcommand{\urlLine}[1]{\newcommand{\makeUrlLine}{#1}}
%%% End Header Definitions %%%

\newcommand{\pr}{\ensuremath{\mathbb{P}}}
\newcommand{\E}{\ensuremath{\mathbb{E}}}


\newcommand{\makeCompact}{%
\setlength{\itemsep}{1pt}%
\setlength{\parsep}{0pt}%
\setlength{\parskip}{0pt}}


\newcommand{\Var}{\mathrm{Var}}


\newcommand{\student}[2]{
\thispagestyle{firststyle}

%%% Begin Header Specification %%%

\headline{Solutions to homework #1}

\makeheader{#2}

\medskip

}

\newcommand{\sol}[1]{ \hspace{.5cm}\\ \noindent \textbf{Exercise #1.}}

% Here you can define your own macros
\newcommand{\calF}{\ensuremath{\mathcal{F}}}
\newcommand{\polyreduce}{$\leq_{m}^{p}$}

\begin{document}
 
 %To create a header for the document use the following macro. \student{Homework number}{Name}

 \student{3}{Stefan Reschke}
 

 \sol{a} \textbf{Proof that \texttt{Max2Sat} is NP-Complete!}

(a) \texttt{Max2Sat} obviously is in NP.

(b) Theorem: \texttt{3Sat} \polyreduce{} \texttt{Max2Sat}, thus \texttt{Max2Sat} is NP-hard.

\vspace{.3cm}
\noindent \textit{Proof.} Given an instance of \texttt{3Sat} $\mathcal{S}_3$, we construct a corresponding instance $\mathcal{S}_{\texttt{max2}}$ of \texttt{Max2Sat}. $\mathcal{S}_3$ has $l$ given clauses $c_1,\ldots,c_l$ consisting of 3 literals. We reconstruct these clauses s.t. we are left with some clauses of length $\leq 2$ that are used in $\mathcal{S}_{\texttt{max2}}$.

Assume we have a clause $c_i=x_i\lor y_i\lor z_i$ of $\mathcal{S}_3$. We add the following clauses to $\mathcal{S}_{\texttt{max2}}$ \begin{gather*}
	\big\{w_i, x_i, y_i, z_i,\\ \lnot x_i\lor \lnot y_i, \lnot y_i\lor \lnot z_i, \lnot x_i\lor \lnot z_i, \\ \lnot w_i\lor x_i, \lnot w_i\lor y_i, \lnot w_i\lor z_i \big\},
\end{gather*} where $w_i$ are newly introduced variables.

If $\mathcal{S}_3$ is satisfiable, at least one of the variables $x_i$, $y_i$ or $z_i$ in a clause $c_i$ must be true. Thus we can simply look at all three cases: \begin{enumerate}
\item If $\texttt{true}, \texttt{false}, \texttt{false}$, then 6 clauses, if $w_i$ is \texttt{true}, or 7 clauses, otherwise, are \texttt{true}.
\item If $\texttt{true}, \texttt{true}, \texttt{false}$, then 7 clauses are \texttt{true}.
\item If $\texttt{true}, \texttt{true}, \texttt{true}$, then 7 clauses, if $w_i$ is \texttt{true}, or 6 clauses, otherwise, are \texttt{true}.
\end{enumerate} So in any case, we can choose $w_i$ such that 7 clauses of these groups are always true. Thus, if we pick $k=7\cdot l$, we should have $\mathcal{S}_{\texttt{max2}}$ set up.

This whole conversion took polynomial time, thus we created $\mathcal{S}_{\texttt{max2}}$ from $\mathcal{S}_{3}$ in polynomial time. \qed
 

 \sol{b} \textbf{Proof that \texttt{MaxCut} is NP-Complete!}

(a) \texttt{MaxCut} obviously is in NP.

(b) Theorem: \texttt{Nae3Sat} \polyreduce{} \texttt{MaxCut}, thus \texttt{MaxCut} is NP-hard.

\vspace{.3cm}
\noindent \textit{Proof.} Given a \texttt{Nae3Sat} instance $\mathcal{S}_{\texttt{nae3}}$, we construct a \texttt{MaxCut} instance $\mathcal{M}$. We assume that $\mathcal{S}_{\texttt{nae3}}$ is given $m$ clauses $c_1,\ldots, c_m$ and contains $n$ variables.
\begin{enumerate}
 \item We construct a graph $G$ such that each variable $x_i$ used in $c_1,\ldots, c_m$ produces vertices $x_i$ and $\lnot x_i$.
 \item We set $10\cdot m$ edges between $x_i$ and $\lnot x_i$.
 \item For every clause $c_i=x_i\lor y_i\lor z_i$, we set edges $(x_i, y_i)$, $(y_i, z_i)$ and $(z_i, x_i)$ in G.
\end{enumerate}
Now we set $k=10mn+2m$, push it all into $\mathcal{M}$ and claim that if $c_1,\ldots c_m$ are satisfiable, $\mathcal{M}$ returns \texttt{true}.

This is easy to see, because either $x_i$ or $\lnot x_i$ hold, thus all edges created in step 2 are cut (which contributes $10mn$ to our goal of at least $k$). Since the variables appear in clauses and are near the cut done in the previous step, we observe that for each clause 2 additional edges go across the cut (which contributes the missing $2m$ edges to the cut). \qed


 \sol{c} \textbf{Proof that \texttt{DomSet} is NP-Complete!}

(a) \texttt{DomSet} obviously is in NP.

(b) Theorem: \texttt{VC} \polyreduce{} \texttt{DomSet}, thus \texttt{DomSet} is NP-hard.

\vspace{.3cm}
\noindent \textit{Proof.} Given an instance of \texttt{VC}, we alter its input graph $G$ and value $k$ to get a corresponding instance of \texttt{DomSet}. \begin{enumerate}
 \item Let $G'$ be a copy of $G$.
 \item For each edge $(x,y)$ in $G'$, we introduce a new vertex $w$.
 \item For each added $w$ we add edges $(x, w)$ and $(w, y)$ in $G'$.
 \item Let $k'$ be $k+n$ where $n$ is the number of isolated vertices in $G$.
\end{enumerate}
We claim that $(G', k')$ is a correctly matched instance of \texttt{DomSet}.

 
\end{document}
